problem:  
Let \((M,\mathcal D,g)\) be a **compact, equiregular contact sub‑Riemannian manifold** with canonical Popp measure \(\mu\).  Denote by \(\mathrm{Ric}_{\mathrm{ABR}}\) the **Agrachev–Barilari–Rizzi curvature** along normal extremals.  Assume that \((M,d_{\mathrm{subR}},\mu)\) satisfies the **Lott–Sturm–Villani curvature–dimension inequality** \(CD(K,N)\) for some finite \(N>2n+1\), interpreted via entropy convexity along Wasserstein geodesics of the sub‑Riemannian squared distance.  

Recall that Juillet and subsequent authors showed that even the Heisenberg group with its standard sub‑Riemannian structure and Haar measure does **not** satisfy any finite‑dimensional \(CD(K,N)\) condition.  Consequently, sub‑Riemannian \(CD(K,N)\) spaces, if they exist, must exhibit exceptional geometric rigidity.  

**Problem:**  
Assume that \((M,\mathcal D,g,\mu)\) is contact and equiregular and that for some \(K>0,N<\infty\), the metric‑measure space \((M,d_{\mathrm{subR}},\mu)\) satisfies \(CD(K,N)\).  
Must \((M,\mathcal D,g)\) then be **locally Riemannian**, i.e. must the underlying contact distribution actually arise from a Riemannian structure (so that \(\mathcal{D}=TM\))?  
Equivalently, do the only contact sub‑Riemannian manifolds supporting a finite‑dimensional \(CD(K,N)\) curvature lower bound coincide (up to local isometry) with Riemannian manifolds?  

More concretely:
1. Can one prove that the \(CD(K,N)\) convexity of entropy along all normal Wasserstein geodesics forces the vertical Reeb direction to be an *integrable flat factor*, thereby eliminating the non‑holonomic structure?  
2. If not, can one classify possible **non‑Riemannian contact models** (if any) capable of supporting true \(CD(K,N)\) bounds despite the presence of nontrivial bracket relations?  

Why is it a "good" problem:  
- **Anchored in known results but truly open:** It builds on Juillet’s and others’ theorems showing failure of \(CD(K,N)\) for the Heisenberg group, and asks whether that failure is universal among contact (hence step‑2) sub‑Riemannian geometries or whether any genuinely non‑Riemannian example can exist.  No such classification is known.  
- **Mathematically precise and falsifiable:** The problem posits a crisp rigidity statement: “A contact sub‑Riemannian manifold satisfying a finite‑dimensional \(CD(K,N)\) must be Riemannian.”  This can be proven true or false; both outcomes lead to deep insight.  
- **Bridges multiple areas:** It unifies the synthetic curvature‑dimension framework of optimal transport, canonical Ricci notions (ABR curvature), and contact/sub‑Riemannian geometry.  Showing such rigidity (or finding counterexamples) would profoundly clarify how transport‑based curvature interacts with nonholonomy.  
- **Extremely simple yet profound:** The question has a minimalistic formulation understandable in one sentence—“Are finite‑dimensional \(CD(K,N)\) spaces within contact geometry necessarily Riemannian?”—but resolving it would pinpoint the precise obstruction that nonholonomy poses to synthetic curvature theory.  
- **Potential impact:** A resolution would either establish an unexpected **synthetic‑curvature rigidity theorem** (implying that the CD framework is essentially Riemannian) or unveil the first genuine **non‑Riemannian CD space**, reshaping our conceptual map of curvature in metric measure geometry.  It exemplifies the “narrow entrance, profound depth” hallmark of a good research problem.